%-----------------------Homework------------------------------------
%-------------------Arman Shokrollahi---------------------------------
%---------------------Coding Theory-------------------------------

\documentclass[a4 paper]{article}
% Set target color model to RGB
\input{preamble.tex}

\begin{document}
\homework{Práctia \#6}{Junio. 05 al 08 de 2019}{MSc. Fausto M. Lagos S.}{}{Aquí va su nombre}{}

En esta plan lector se desarrollarán dos planteamientos el primero correspondiente a tres puntos y el segundo a cuatro puntos. El objetivo es afianzar sus conocimientos antes de la Prueba de Evaluación Continua correspondiente, puede utilizar todo el material bibliográfica a su disposición y también preguntar todo lo que considere necesario. Preste atención al margen derecho donde encontrará premios y bonificaciones adicionales.

\vspace{5mm}

\problem{1} \textbf{4 puntos} - Del Cálculo de Spivak resuelva el enunciado (iii) edel numeral 1 de la página 284, apoyese en la lectura de la sección II - Significado de la derivada -.
\vskip 3mm
\solution{}
\vspace{30mm}

\problem{2} \textbf{3 puntos} - Construya la gráfica de la función del problema 1 y muestre los puntos máximos y mínimos locales.
\vskip 3mm
\solution{}
\end{document}